%! Bias in Samples and Surveys — Unit 1, Lesson 1.4 — Group Activity
\documentclass[10pt]{article}
\usepackage{saar-article}
\usepackage{saar-boxes}
\newcommand{\TallMath}[1]{$\displaystyle #1\rule[-1.4em]{0pt}{3.2em}$}

\begin{document}

\pageheader{Unit 1, Lesson 1.4}{Group Activity}
% NAMESTRIP: no \namedateperiod here. The student writes their name once, on the
% cover this component is stapled behind. See references/conventions.md.

\vspace{0.15in}

\begin{headlinebox}{frost}
\small
\textbf{The Cards That Came Back --- Cedar Ridge Apartments.} Cedar Ridge has
$600$ households in $30$ buildings of $20$. Management mails a parking survey to
\emph{every} household. $150$ cards come back, and $96$ of those say parking is a
problem. Management is about to tell the owners that Cedar Ridge needs a bigger
lot. \emph{(Recall the layout from last class: buildings $1$--$6$ are studios,
$7$--$20$ are one-bedrooms, $21$--$30$ are two-bedrooms.)}
\end{headlinebox}

\vspace{0.06in}

\boxguard[30]
\begin{tcolorbox}[colback=white, colframe=black!40, breakable,
    title=\textbf{Tier R --- Remediate},
    fonttitle=\bfseries, arc=2mm, left=3mm, right=3mm, top=1.5mm, bottom=1.5mm]
\small
\begin{enumerate}[leftmargin=1.6em, itemsep=4pt, topsep=3pt]

  \item Name the two numbers in this study.

        Population: \blank{5cm} \hfill Households that answered: \blank{1.8cm}

  \item Name the main source of bias in each plan. Choose from
        \emph{undercoverage}, \emph{nonresponse}, \emph{social desirability},
        \emph{demand bias}, and \emph{question order bias}.

        \begin{center}
        \renewcommand{\arraystretch}{1.4}
        \begin{tabularx}{0.96\linewidth}{@{} X p{3.2cm} @{}}
        \toprule
        \textbf{The plan} & \textbf{Source of bias} \\
        \midrule
        Mail a card to all $600$ households; $150$ mail it back.
          & \blank{2.8cm} \\
        Survey only the households that bought a parking sticker.
          & \blank{2.8cm} \\
        The building manager stands beside the resident while they fill it in.
          & \blank{2.8cm} \\
        The card is headed \emph{``Help us prove we need more parking.''}
          & \blank{2.8cm} \\
        ``How crowded is the lot?'' is asked \emph{before} ``Should we spend
        \$60,000 on more spaces?'' & \blank{2.8cm} \\
        \bottomrule
        \end{tabularx}
        \end{center}

  \item What percent of the $600$ households answered at all?

        \begin{work}
          \dfrac{150}{600} &= 0.25 \\
                           &= 25\%
        \end{work}

  \item Of the $150$ households that answered, $96$ said parking is a problem.
        What percent of the answering households is that?

        \begin{work}
          \dfrac{96}{150} &= 0.64 \\
                          &= 64\%
        \end{work}

  \item One plan in item 2 reached every household on the frame and still went
        wrong. Which one, and what went wrong?
        \writeline
\end{enumerate}
\end{tcolorbox}

\vspace{0.06in}

\boxguard
\begin{tcolorbox}[colback=white, colframe=black!40, breakable,
    title=\textbf{Tier A --- Approaching Proficiency},
    fonttitle=\bfseries, arc=2mm, left=3mm, right=3mm, top=2mm, bottom=2mm]
\small
\begin{enumerate}[leftmargin=1.6em, itemsep=5pt, topsep=3pt]

  \item Management uses its $64\%$ to describe all $600$ households.

        \textbf{(a)} How many of the $600$ households does $64\%$ predict?

        \begin{work}
          0.64 \times 600 &= 384 \text{ households}
        \end{work}

        \textbf{(b)} A year later the office asks \emph{every} household and
        finds that $240$ of the $600$ say parking is a problem. What percent is
        that?

        \begin{work}
          \dfrac{240}{600} &= 0.40 \\
                           &= 40\%
        \end{work}

        \textbf{(c)} The mail-in study overestimated by \blank{2cm} percentage
        points, which is \blank{2cm} households too many.

  \item Management says the first study just needed more cards. They mail again,
        get $300$ back this time, and $192$ of those say parking is a problem.

        \textbf{(a)} What percent is that?

        \begin{work}
          \dfrac{192}{300} &= 0.64 \\
                           &= 64\%
        \end{work}

        \textbf{(b)} Twice as many cards, same answer, still $24$ points above
        the truth. Explain why the extra cards did not help.
        \writelines{2}

  \item The card asks: \emph{``Don't you agree that Cedar Ridge's parking lot is
        dangerously overcrowded?''}

        \textbf{(a)} Name the kind of response bias this wording invites.
        \blank{4.5cm}

        \textbf{(b)} Rewrite the question so it does not push either way.
        \writelines{2}
\end{enumerate}
\end{tcolorbox}

\vspace{0.06in}

\boxguard
\begin{tcolorbox}[colback=white, colframe=black!40, breakable,
    title=\textbf{Tier E --- Extension},
    fonttitle=\bfseries, arc=2mm, left=3mm, right=3mm, top=2mm, bottom=2mm]
\small
\begin{enumerate}[leftmargin=1.6em, itemsep=5pt, topsep=3pt]

  \item Every household got a card, so nobody was left off the frame. Explain why
        the $450$ households that threw theirs away are still a problem --- and
        say which direction they push the $64\%$.
        \writelines{2}

  \item Two-bedroom households own more cars than studio households do. Suppose
        two-bedroom households mailed their cards back at twice the rate of
        studio households.

        \textbf{(a)} Whose opinions are overrepresented among the $150$ cards?
        \writeline

        \textbf{(b)} Does that make the $64\%$ too high or too low? Explain how
        you know.
        \writelines{2}

  \item Management can afford exactly \emph{one} change. For each option, name
        the bias it removes and what it costs them.

        \textbf{(a)} Knock on the $450$ silent doors.
        \writeline

        \textbf{(b)} Interview a stratified random sample of $60$.
        \writeline

        \textbf{(c)} Reword the card; let residents return it anonymously.
        \writeline
\end{enumerate}
\end{tcolorbox}

\end{document}
